\chapter{Versioning and releases}

Release automation is only reliable if your versioning strategy is clear.
Decide how versions are produced and how artifacts are published.

\section{Semantic versioning}
Use \code{MAJOR.MINOR.PATCH} where:
\begin{itemize}[nosep]
  \item MAJOR breaks compatibility
  \item MINOR adds features
  \item PATCH fixes bugs
\end{itemize}

\section{Tag-driven releases}
Tagging a commit is the simplest way to trigger a release pipeline.

\begin{lstlisting}[language=YAML]
name: Release

on:
  push:
    tags:
      - "v*.*.*"

permissions:
  contents: write

jobs:
  build:
    runs-on: ubuntu-latest
    steps:
      - uses: actions/checkout@v4
      - run: ./scripts/build
      - uses: actions/upload-artifact@v4
        with:
          name: dist
          path: dist/
\end{lstlisting}

\section{Release branches}
For long-lived products, you may need maintenance branches.
Keep release workflows identical across branches and use tags to signal releases.

\section{Changelogs}
Generate changelogs as part of the release.
Use commit conventions or labels to make automation reliable.

\section{Manual releases}
A manual release is still automation.
Use \code{workflow\_dispatch} inputs to select version, environment, and dry-run.

\begin{lstlisting}[language=YAML]
on:
  workflow_dispatch:
    inputs:
      version:
        description: "Release version"
        required: true
      dry_run:
        description: "Skip publish"
        required: false
        default: "true"
\end{lstlisting}

\section{Release promotion}
Promote the same artifact from staging to production instead of rebuilding.
This keeps releases reproducible and easier to audit.

\section{Signing and provenance}
If you ship binaries, consider signing them.
Use attestations or provenance to prove where and how the build was produced.

\begin{patternbox}{Pattern: release is a workflow, not a human checklist}
If you can ship by running \code{workflow\_dispatch} and selecting a version,
you can ship on Friday with confidence.
\end{patternbox}

\begin{checklistbox}{Checklist: releases}
\begin{itemize}[nosep]
  \item Tag releases with a consistent prefix.
  \item Store build outputs as artifacts.
  \item Keep release scripts in \code{scripts/}.
\end{itemize}
\end{checklistbox}
